\begin{frame}{$\varepsilon$ 는 ``상한선''인가, ``속도''인가}
  ``110개 중 1개''라는 소박한 변화지만, 이 숫자가
  \kq{$\varepsilon$ 이 ``학습의 \textbf{상한선}''을 제약하는가,
  아니면 단순히 ``\textbf{수렴 속도와 노이즈 크기}''만 조정하는
  가?} 를 생각하게 해준다.
  \vskip0.5em
  \textbf{정답은 전자(대부분)}. $\varepsilon > 0$ 인 한 행동
  정책은 무작위 행동을 섞으므로, 그 정책의 가치 $Q^\pi$ 는
  원래 $Q^{*}$ 보다 낮다. $\varepsilon$ 을 0에 가깝게 하면
  $\pi \to$ 그리디 정책에 가까워져 $Q^\pi \to Q^{*}$ 가
  된다.
  \vskip0.5em
  \begin{block}{$\varepsilon$ 를 0으로 \textbf{완전히} 줄이면
  어떨까}
    그때는 더 이상 새 (상태, 행동) 쌍을 방문하지 못해, 아직
    제대로 추정되지 않은 상태의 $Q$ 가 \textbf{영원히} 그
    자리에서 멈춘다 --- 탐험의 대가(낮은 즉석 보상)를 치르지
    않으면 새로운 지식을 얻을 기회도 사라진다.
    \textbf{``탐험 vs 착취'' 트레이드오프의 가장 정직한 모습.}
  \end{block}
\end{frame}
